\section{干预对比图谱}

\begin{itemize}[nosep]
  \item \textbf{窗口捕获。} 目标：时机与进入杠杆。证据车道：strong + design inference。用途：把尝试移动到一个边界上，例如起床、休息后重新进入，或通勤转变时刻。
  \item \textbf{状态预载。} 目标：切换成本与路径清晰度。证据车道：medium + design inference。用途：在窗口打开前，预先准备好准确的下一页、下一句、下一条路线或下一件工具。
  \item \textbf{环境再加权。} 目标：线索景观与局部强化。证据车道：strong/medium + design inference。用途：移除那些会重建旧状态的线索，并加入让目标状态更容易进入的线索。
  \item \textbf{历史塑形。} 目标：涌现与路径依赖。证据车道：medium + design inference。用途：在相同情境中反复使用同一个进入动作，让局部历史积累起来。
  \item \textbf{最小可行进入。} 目标：在有限意志条件下完成障碍跨越。证据车道：medium + design inference。用途：定义一个仍然算进入目标机制的最小真实动作。
  \item \textbf{外部脚手架。} 目标：有效控制输入。证据车道：medium + design inference。用途：当私人控制过于嘈杂，难以可靠触发进入时，使用时间表、伙伴或截止期。
\end{itemize}